\chapter*{Preface}
\addcontentsline{toc}{chapter}{Preface}
\markboth{Preface}{Preface}

By the time you reach category theory, you've seen the same
pattern --- a structure, the maps between structures that preserve it,
and a notion of when two things are ``the same'' relative to those maps
--- repeat across groups, vector spaces, topological spaces, and more.
This booklet makes that pattern the actual subject: categories,
functors and natural transformations, adjunctions, limits and colimits,
monads, universal algebra, representation theory, abelian categories
and homological algebra, derived functors, topos theory, and a first
look at $\infty$-categories.

This is deliberately the most demanding booklet to place in the
sequence, because it depends a little on almost everything before it
--- not as detailed prerequisite knowledge, but as the set of examples
that make the abstraction worth doing at all.
